\documentclass[12pt,a4paper]{ctexart}
% 公共排版文件（使用 XeLaTeX 编译）
\usepackage[a4paper,margin=2.35cm]{geometry}
\usepackage{booktabs,longtable,tabularx,array,multirow,enumitem}
\usepackage{graphicx,xcolor,amsmath,hyperref,lastpage,fancyhdr}
\usepackage{tikz}
\usetikzlibrary{arrows.meta,positioning,shapes.geometric,fit,calc}
\hypersetup{colorlinks=true,linkcolor=blue!60!black,urlcolor=blue!60!black,citecolor=blue!60!black,pdftitle={逸仙活动云课程报告}}
\setlist{itemsep=2pt,topsep=3pt,leftmargin=2em}
\setlength{\parindent}{2em}
\setlength{\headheight}{15pt}
\renewcommand{\arraystretch}{1.35}
\newcolumntype{Y}{>{\raggedright\arraybackslash}X}
\newcommand{\reportcover}[1]{%
  \begin{titlepage}
  \begin{center}
  \vspace*{3cm}{\zihao{1}\bfseries 中山大学\\[0.6cm]《软件工程》期末大作业\par}
  \vspace{2.5cm}{\zihao{2}\bfseries #1\par}
  \end{center}
  \vspace{2cm}
  \begin{center}
  \begin{tabular}{l}
  {\Large 项目名称：逸仙活动云——中山大学校园活动云平台}\\
  {\Large 小组成员：钟梓俊\quad 邱剑波\quad 张宸睿}\\
  {\Large 指导老师：王可泽}
  \end{tabular}
  \end{center}
  \vfill
  \begin{center}
  {\large 2026 年 7 月 16 日\par}
  \end{center}
  \end{titlepage}}

\pagestyle{fancy}\fancyhf{}\lhead{逸仙活动云}\rhead{软件工程期末大作业}\cfoot{第 \thepage\ 页，共 \pageref{LastPage} 页}
\setcounter{tocdepth}{2}

\begin{document}
\reportcover{软件配置与运维文档}
\pagenumbering{Roman}\tableofcontents\clearpage\pagenumbering{arabic}

\section{文档概述}
本文档定义逸仙活动云项目的配置管理、版本管理、持续集成、软件部署、运行维护及应急方案。目标是确保每个交付版本可识别、可重建、可部署、可监控和可恢复。文档适用于项目团队的开发和运维人员，以及后续接手项目维护的第三方开发者。文档中“已实现”指代码库已有对应实现；“计划”指建议在交付或运维阶段落实，二者不作混同。

\section{配置管理}
\subsection{配置项说明}
项目涉及的配置项按类别划分如下：
\begin{enumerate}
\item \textbf{应用配置：} 前端 API 地址（\texttt{VITE\_API\_BASE\_URL}）、后端数据库连接串（\texttt{DATABASE\_URL}）、Redis 连接串（\texttt{REDIS\_URL}）、CORS 域名。
\item \textbf{安全配置：} JWT 签名密钥（\texttt{JWT\_SECRET\_KEY}）、密码哈希算法、验证码有效期。
\item \textbf{外部服务配置：} 邮件 SMTP 服务器地址与凭据、外部搜索 API 密钥、AI 服务 API 端点和密钥、数据源抓取间隔与超时。
\item \textbf{部署配置：} 容器镜像版本、数据卷挂载路径、Gunicorn Worker 数量、健康检查端点。
\end{enumerate}

\subsection{配置文件管理}
项目通过分层配置文件管理各类配置：
\begin{enumerate}
\item \textbf{前端：} \texttt{.env.example} 作为模板文件提交到仓库，开发者复制为 \texttt{.env.local} 后填充实际值。Vite 在开发阶段自动加载根目录的 \texttt{.env.local} 和环境变量。
\item \textbf{后端：} \texttt{.env.example} 作为模板提交到仓库，开发者复制为 \texttt{.env} 后填充。Flask 通过 \texttt{python-dotenv} 在应用启动时加载。
\item \textbf{Docker Compose：} \texttt{docker-compose.yml} 定义服务拓扑和卷挂载，通过 \texttt{environment} 字段或 \texttt{env\_file} 引用环境变量文件。
\end{enumerate}
所有配置文件模板不包含真实密钥和密码。开发环境、测试环境和生产环境使用各自独立的环境文件，互不干扰。

\subsection{环境变量管理}
\begin{center}
\begin{tabularx}{0.9\textwidth}{p{2.8cm}p{5.5cm}p{5.5cm}}
\toprule 环境 & 配置来源 & 数据与依赖 \\\midrule
本地开发 & 项目根目录 \texttt{.env.local}（前端）或 \texttt{.env}（后端） & Mock 服务或 SQLite 数据库；不连接生产 Redis 或邮件服务。 \\
测试/CI & GitHub Actions Secret 或临时环境变量 & 内存 SQLite；禁用真实邮件、Redis 和外部 LLM。 \\
预发布 & 独立 \texttt{.env} 文件或容器编排工具注入 & 独立 PostgreSQL 和 Redis 实例；脱敏样本数据。 \\
生产 & 受控 Secret 管理服务（如 Docker Secret 或云平台密钥管理） & PostgreSQL 主从、Redis 集群；HTTPS、备份和监控必须启用。 \\
\bottomrule
\end{tabularx}
\end{center}

\subsection{敏感信息管理}
\begin{enumerate}
\item \textbf{密钥分类：} JWT 签名密钥、邮件 SMTP 密码、AI/搜索 API 密钥、数据库连接密码。
\item \textbf{存储原则：} 敏感信息以环境变量或 Secret 管理服务注入容器，不写入代码仓库、日志文件或截图。
\item \textbf{泄露处置：} 发现泄露后立即撤销原密钥、生成新密钥、轮换受影响范围内的所有凭据、审计泄露时间窗口内的访问记录、重建受影响的制品。
\item \textbf{开发默认值：} \texttt{.env.example} 中的默认凭据仅用于本地演示环境，生产部署前必须替换为高熵随机值。
\end{enumerate}

\section{版本管理}
\subsection{版本号规范}
项目使用语义化版本号 \texttt{vMAJOR.MINOR.PATCH}：
\begin{enumerate}
\item \textbf{MAJOR：} 破坏性接口变更或数据库 Schema 不向后兼容时递增。
\item \textbf{MINOR：} 新增向后兼容的功能时递增。
\item \textbf{PATCH：} 缺陷修复或内部优化时递增。
\end{enumerate}
前端和后端仓库独立标识版本。发布时以兼容版本对组成系统基线，同时记录前端 tag、后端 tag、接口契约版本和数据库 Schema 版本。

\subsection{发布版本管理}
\begin{enumerate}
\item 主分支（main）始终保持可构建状态，所有变更通过合并请求合入。
\item 功能分支在短生命周期内开发，完成后删除。
\item 每个可交付版本从 main 创建 tag，tag 名使用 \texttt{vMAJOR.MINOR.PATCH}。
\item 紧急修复从受影响 tag 创建 hotfix 分支，修复后同时合回 main 和发布分支。
\end{enumerate}

\subsection{变更记录}
所有版本变更记录在 release notes 中，包含以下信息：
\begin{enumerate}
\item 版本号和发布日期。
\item 新增功能清单及关联需求编号（FR-xx）。
\item 修复的缺陷清单及关联缺陷编号。
\item 接口或数据库 Schema 变更摘要。
\item 升级注意事项（如有不兼容变更）。
\end{enumerate}

\subsection{回滚机制}
\begin{enumerate}
\item \textbf{Docker 部署回滚：} 停止当前容器，使用上一版本的镜像重新启动，恢复上一版本的数据库备份。
\item \textbf{数据库回滚：} 在迁移前执行数据库全量备份；回滚时恢复备份并验证数据完整性。
\item \textbf{前端回滚：} 将反向代理或 Web 服务器的静态资源目录切换回上一版本的构建产物。
\end{enumerate}
回滚后保留当前版本和回滚版本的日志、镜像摘要和配置快照，用于后续根因分析。

\section{持续集成}
\subsection{CI 流程}
前端仓库已配置 GitHub Actions CI 工作流（\texttt{.github/workflows/ci.yml}），在 push 和 pull request 事件上触发。工作流执行以下步骤：
\begin{enumerate}
\item 检出代码并配置 Node 22 环境（\texttt{actions/setup-node@v4}）。
\item 执行 \texttt{npm ci} 安装精确版本依赖。
\item 安装 Chromium 浏览器（\texttt{npx playwright install --with-deps chromium}）。
\item 执行 \texttt{npm run build}（类型检查 + 生产构建）。
\item 执行 \texttt{npm run test:unit}（Vitest 单元测试）。
\item 执行 \texttt{npm run test:e2e}（Playwright 端到端测试）。
\end{enumerate}

\subsection{自动化构建}
前端通过 \texttt{npm run build} 触发 \texttt{vue-tsc -b \&\& vite build}，串行执 TypeScript 类型检查和 Vite 生产构建，输出可部署的静态资源到 \texttt{dist/} 目录。后端通过 Dockerfile 定义构建过程，安装 Python 依赖并打包 Flask 应用。构建失败时 CI 工作流终止并报告错误。

\subsection{自动化测试}
\begin{enumerate}
\item \textbf{单元测试：} Vitest（前端）覆盖组件渲染、状态管理和工具函数；pytest（后端）覆盖 API、服务和业务规则。
\item \textbf{端到端测试：} Playwright 覆盖访客浏览、用户注册登录、活动搜索等关键用户旅程。
\item \textbf{聚合检查：} \texttt{npm run verify} 串行执行构建、单元测试和端到端测试，作为提交前的完整性验证。
\end{enumerate}

\subsection{构建产物管理}
前端构建产物为静态文件（HTML/CSS/JS），直接部署到 Web 服务器或反向代理。后端构建产物为 Docker 镜像，推送到镜像仓库后由 Compose 引用。发布制品包括前端静态构建包、后端不可变镜像、SBOM 依赖清单、测试报告和 release notes，以 tag 和 commit SHA 双重标识。

\section{软件部署}
\subsection{环境要求}
\noindent\begin{tabularx}{\textwidth}{p{2.5cm}Y}
\toprule 项目 & 要求 \\\midrule
操作系统 & Linux（推荐 Ubuntu 22.04+）/ Windows / macOS（仅开发） \\
运行环境 & Docker Engine 24+、Docker Compose v2 \\
数据库 & PostgreSQL 15+（部署）/ SQLite 3（开发） \\
缓存 & Redis 7+（可选，部署推荐启用） \\
反向代理 & Nginx（生产，建议配置 TLS） \\
网络 & 容器间内网通信；生产环境需开放 443（HTTPS） \\
\bottomrule
\end{tabularx}

\subsection{安装部署步骤}
以下步骤涉及前端和后端两个仓库，请根据提示在对应目录下执行：
\begin{enumerate}
\item \textbf{后端目录：} 复制环境变量模板并填写配置：\texttt{cp .env.example .env}。
\item \textbf{后端目录：} 执行 \texttt{docker compose up -d --build} 启动所有服务（API、PostgreSQL、Redis、Celery Worker/Beat）。
\item 访问 \texttt{GET http://127.0.0.1:5000/api/health} 确认 API 存活。
\item \textbf{前端目录：} 执行 \texttt{npm run build} 构建前端静态资源。
\item 将前端构建产物 \texttt{dist/} 目录部署到 Nginx 或反向代理的静态资源路径。
\item 验证登录、活动查询和审核等最小业务流程正常。
\end{enumerate}

\subsection{初始化配置}
首次启动时，Flask 应用通过 \texttt{app.config[\"SQLALCHEMY\_TRACK\_MODIFICATIONS\"]} 和模型注册自动创建数据库表。此方式仅适用于开发和演示环境，生产环境必须使用受控迁移工具（如 Alembic）管理 Schema 变更。系统同时提供开发用初始数据（默认管理员账号和示例活动），首次部署后应修改默认密码。

\subsection{升级步骤}
\begin{enumerate}
\item 拉取目标版本的 Docker 镜像或前端构建产物。
\item 备份当前数据库（\texttt{pg\_dump}）和配置文件。
\item 执行数据库迁移脚本（如有 Schema 变更）。
\item 更新 Docker 镜像版本标签并重启服务：\texttt{docker compose up -d --pull always}。
\item 执行冒烟测试验证升级后业务正常。
\end{enumerate}

\subsection{卸载步骤}
\begin{enumerate}
\item 停止所有容器：\texttt{docker compose down}。
\item 删除数据卷（谨慎操作，将丢失所有数据）：\texttt{docker volume rm <volume\_name>}。
\item 删除镜像：\texttt{docker rmi <image\_name>:<tag>}。
\item 删除项目目录和配置文件。
\end{enumerate}

\section{运行维护}
\subsection{日志管理}
\begin{enumerate}
\item \textbf{日志来源：} Flask 应用日志（请求、错误、审计）、Celery Worker 日志（任务执行）、Nginx 访问日志、PostgreSQL 慢查询日志。
\item \textbf{日志内容规范：} 日志不得包含密码、JWT、完整邮箱验证码或完整搜索词。审核和管理操作必须记录操作者、动作、目标和时间。
\item \textbf{日志轮转：} 使用 Docker 的日志驱动或 logrotate 工具实现日志轮转，避免磁盘空间耗尽。
\end{enumerate}

\subsection{数据备份}
\begin{enumerate}
\item \textbf{备份策略：} PostgreSQL 数据库执行每日全量备份，结合 WAL 日志实现增量恢复。
\item \textbf{备份存储：} 备份文件加密后存储于异地或独立存储系统，与生产环境隔离。
\item \textbf{验证：} 定期执行恢复演练，确保备份文件可正常恢复且数据完整性通过校验。
\end{enumerate}

\subsection{数据恢复}
\begin{enumerate}
\item 确认故障范围和影响程度，隔离故障环境。
\item 选择最近的可用备份文件，执行恢复操作。
\item 恢复后验证核心数据完整性（用户数、活动数、关系数据）。
\item 确认服务健康检查通过后，重新开放服务。
\end{enumerate}

\subsection{异常处理}
\noindent\begin{tabularx}{\textwidth}{p{2.8cm}Y p{3.5cm}}
\toprule 异常类型 & 处理措施 & 告警方式 \\\midrule
API 服务不可用 & 检查容器状态、应用日志和数据库连接；重启服务或回滚到上一版本。 & 健康检查连续失败触发告警。 \\
数据库连接失败 & 检查 PostgreSQL 容器、数据卷和网络；从备份恢复数据卷。 & 应用日志记录连接错误并告警。 \\
Redis 不可达 & 确认 Redis 容器运行状态；限流和队列功能降级，核心 API 继续运行。 & 日志记录降级状态，待连接恢复后自动重连。 \\
任务队列积压 & 检查 Worker 进程和外部依赖；暂停失败率高的数据源。 & 监控队列深度，超过阈值触发告警。 \\
安全事件 & 撤销泄露的凭据，保全日志，评估数据影响范围，修复漏洞。 & 异常登录或越权行为记录审计日志。 \\
\bottomrule
\end{tabularx}

\subsection{故障排查}
\begin{enumerate}
\item \textbf{容器级别：} 使用 \texttt{docker ps} 检查容器运行状态，\texttt{docker logs <container>} 查看应用日志。
\item \textbf{应用级别：} 访问 \texttt{/api/health} 检查 API 存活状态；查看 Flask 日志定位 5xx 错误；检查数据库连接池和慢查询。
\item \textbf{网络级别：} 使用 \texttt{docker network inspect} 检查容器间通信；确认反向代理配置和防火墙规则。
\item \textbf{数据级别：} 检查数据库表结构完整性；确认外键约束和唯一约束未损坏；检查数据卷挂载状态。
\end{enumerate}

\section{发布与回滚}
发布流程遵循以下步骤：
\begin{center}
\begin{tikzpicture}[font=\small,
  node distance=0.55cm and 1.1cm,
  flowstep/.style={draw,rounded corners,align=center,minimum width=2.35cm,minimum height=.7cm,fill=blue!7},
  arr/.style={-Latex,thick}
]
\node[flowstep] (a) {1. 选择已验证 tag\\冻结变更};
\node[flowstep] (b) [right=of a] {2. 构建/签名制品\\记录摘要};
\node[flowstep] (c) [right=of b] {3. 备份数据库\\检查迁移};
\node[flowstep] (d) [below=of c] {4. 部署 Compose/服务\\注入 Secret};
\node[flowstep] (e) [left=of d] {5. 健康与冒烟\\日志/指标检查};
\node[flowstep] (f) [left=of e] {6. 宣布或回滚\\归档证据};
\draw[arr] (a)--(b);  \draw[arr] (b)--(c);
\draw[arr] (c)--(d);  \draw[arr] (d)--(e);
\draw[arr] (e)--(f);
\end{tikzpicture}
\vspace{0.2em}
{\par\small\itshape 图1：发布与回滚流程}
\end{center}

发布前检查清单：
\begin{enumerate}
\item 选择已通过质量门禁的前后端 commit/tag，冻结接口和数据库迁移版本。
\item 确认无真实密钥、\texttt{.env} 文件或本地数据进入制品。
\item 在预发布环境恢复备份、执行迁移和冒烟测试。
\item 复核 TLS 配置、CORS 设置、请求体限制、健康检查和日志脱敏。
\item 部署后验证健康接口、登录、活动浏览、发布审核和审计查询。
\end{enumerate}
回滚触发条件：健康检查连续失败、核心冒烟测试未通过、错误指标超过阈值。回滚时停止当前服务，恢复上一版本镜像和数据库备份，保留日志和版本信息用于根因分析。

\section{运维风险与应急方案}
\noindent\begin{tabularx}{\textwidth}{p{3cm}Y p{4.5cm}}
\toprule 风险 & 影响 & 应急方案 \\\midrule
数据库损坏或数据丢失 & 用户数据不可恢复，核心业务中断 & 从每日备份恢复；启用只读维护页面；排查损坏原因。 \\
密钥泄露 & 未授权访问、数据泄露 & 立即撤销并轮换密钥；审计泄露时间窗口内操作；通知受影响用户。 \\
外部服务不可用 & 邮件发送、搜索、AI 功能降级 & 切换至降级模式；保留错误日志待服务恢复后重试。 \\
容器运行时故障 & 服务中断 & 检查 Docker 守护进程状态；重新启动容器；切换至备用节点。 \\
\bottomrule
\end{tabularx}

在课程项目范围内，本文档交付了以下可复现的运维资产：
\begin{enumerate}
\item \textbf{容器化部署拓扑：} Dockerfile + docker-compose.yml 定义 API、PostgreSQL、Redis、Celery Worker/Beat 的完整服务编排，单命令（\texttt{docker compose up -d --build}）即可在任意 Linux 服务器上重建运行环境。
\item \textbf{健康检查与日志：} \texttt{/api/health} 端点提供存活检测，Flask 应用日志和 Celery Worker 日志记录运行状态。
\item \textbf{配置管理框架：} \texttt{.env.example} 模板管理环境变量，密钥注入使用环境变量而非硬编码，支持开发/测试/预发布/生产四环境隔离。
\item \textbf{运维操作文档：} 本文档涵盖配置管理、版本管理、CI/CD、部署步骤、升级/卸载、运行维护、异常处理和故障排查等完整操作流程。
\end{enumerate}
生产环境所需的 TLS、集中监控、自动备份、分支保护和灾备演练属于生产加固（production hardening）范畴，与课程项目交付范围相分离。上述加固措施已在本文档中明确记录操作步骤和配置方法，上线前可据此实施。

每次运行事故由值班人员先止损和保全日志，再由开发、测试、运维人员共同完成根因分析、修复、回归验证和预防行动跟踪。

\end{document}
